% Worked Examples: Concept 3.9.2 - Stability and Performance in the Presence of Uncertainty
\documentclass[11pt,a4paper]{article}
\usepackage{worked-examples-template}

\title{\textbf{Worked Examples}\\[0.5em]
\large Concept 3.9.2: Stability and Performance in the Presence of Uncertainty\\[0.3em]
\normalsize \coursesubtitle{} -- \coursetitle}
\author{Assoc.\ Prof.\ William Robertson \and Dr Sean McGowan}
\date{\today}

\begin{document}

\maketitle
\tableofcontents
\newpage

\section{Concept 3.9.2: Robust Stability}

\subsection{Example 1: Stability Bounds with Multiplicative Perturbations}

Robust stability ensures a control system remains stable despite model uncertainties. This example applies the multiplicative uncertainty condition.

\paragraph{Problem:}
Consider a feedback system with nominal process $P(s) = 5/(s+2)$ and controller $C(s) = 2$. 

(a) Calculate the complementary sensitivity function $T(s)$.

(b) Determine the maximum magnitude of $T$ over all frequencies: $M_t = \sup_\omega |T(\ii\omega)|$.

(c) Using the robust stability condition $|\delta(\ii\omega)| < 1/|T(\ii\omega)|$, find the maximum allowable multiplicative uncertainty at $\omega = 0$, $\omega = 2$, and $\omega = 10$ rad/s.

(d) Interpret the results for controller design.

\paragraph{Solution:}

\textbf{Step 1: Loop transfer function and complementary sensitivity}

The loop transfer function is:
\[
L(s) = P(s)C(s) = \frac{5}{s+2} \cdot 2 = \frac{10}{s+2}
\]

The complementary sensitivity function is:
\[
T(s) = \frac{L(s)}{1 + L(s)} = \frac{10/(s+2)}{1 + 10/(s+2)} = \frac{10}{s+2+10} = \frac{10}{s+12}
\]

\textbf{Step 2: Frequency response of $T$}

At frequency $\omega$, substitute $s = \ii\omega$:
\[
T(\ii\omega) = \frac{10}{\ii\omega + 12}
\]

The magnitude is:
\[
|T(\ii\omega)| = \frac{10}{\sqrt{\omega^2 + 144}}
\]

\textbf{Step 3: Find maximum magnitude $M_t$}

Since $|T(\ii\omega)|$ is monotonically decreasing with $\omega$:
\[
M_t = \max_\omega |T(\ii\omega)| = |T(\ii \cdot 0)| = \frac{10}{12} = 0.833
\]

This occurs at $\omega = 0$ (DC gain).

\textbf{Step 4: Robust stability bounds at specific frequencies}

The robust stability condition states: $|\delta(\ii\omega)| < 1/|T(\ii\omega)|$

At $\omega = 0$ rad/s:
\[
|T(\ii \cdot 0)| = \frac{10}{12} = 0.833 \quad \Rightarrow \quad |\delta| < \frac{1}{0.833} = 1.20
\]

At $\omega = 2$ rad/s:
\[
|T(\ii \cdot 2)| = \frac{10}{\sqrt{4 + 144}} = \frac{10}{\sqrt{148}} = 0.822 \quad \Rightarrow \quad |\delta| < 1.22
\]

At $\omega = 10$ rad/s:
\[
|T(\ii \cdot 10)| = \frac{10}{\sqrt{100 + 144}} = \frac{10}{\sqrt{244}} = 0.640 \quad \Rightarrow \quad |\delta| < 1.56
\]

\textbf{Step 5: Interpretation}

\begin{itemize}
\item At low frequencies ($\omega \approx 0$), the uncertainty bound is tighter: $|\delta| < 1.20$
\item At high frequencies ($\omega = 10$), larger uncertainties are tolerable: $|\delta| < 1.56$
\item This makes sense: the controller has more gain at low frequencies, making the system more sensitive to model errors
\item The maximum uncertainty across all frequencies is limited by the minimum of $1/|T(\ii\omega)|$, which equals $1/M_t = 1.20$
\end{itemize}

\textbf{Key insights:}
\begin{itemize}
\item Keeping $M_t$ small (equivalently, the sensitivity peak $M_s$ small) improves robustness
\item The complementary sensitivity $T$ determines tolerance to high-frequency uncertainty
\item Trade-off: high controller gain improves disturbance rejection but reduces robustness margins
\end{itemize}

\subsection{Example 2: Performance Under Uncertainty - Load Disturbance Attenuation}

Understanding how uncertainty affects disturbance rejection is essential for robust performance analysis.

\paragraph{Problem:}
A process has nominal model $P_0(s) = 1/(s+1)$ with multiplicative uncertainty $\delta$ such that the actual process is $P(s) = P_0(s)(1+\delta(s))$. A controller $C(s) = 5$ is used.

(a) Calculate the sensitivity function $S(s)$ for the nominal system.

(b) Determine the transfer function $G_{yv}(s)$ from load disturbance $v$ to output $y$ for the nominal system.

(c) If the actual process has $\delta(s) = 0.2$, find the actual $G_{yv}$ and compare with the nominal case.

(d) Calculate the relative error in disturbance rejection performance.

\paragraph{Solution:}

\textbf{Step 1: Nominal sensitivity function}

With nominal process $P_0(s) = 1/(s+1)$ and controller $C(s) = 5$:
\[
L_0(s) = P_0(s)C(s) = \frac{5}{s+1}
\]

The sensitivity function is:
\[
S_0(s) = \frac{1}{1 + L_0(s)} = \frac{1}{1 + 5/(s+1)} = \frac{s+1}{s+1+5} = \frac{s+1}{s+6}
\]

\textbf{Step 2: Load disturbance transfer function (nominal)}

The transfer function from load disturbance to output is:
\[
G_{yv,0}(s) = P_0(s)S_0(s) = \frac{1}{s+1} \cdot \frac{s+1}{s+6} = \frac{1}{s+6}
\]

\textbf{Step 3: Actual system with uncertainty}

With $\delta = 0.2$, the actual process is:
\[
P(s) = P_0(s)(1 + \delta) = \frac{1}{s+1} \cdot 1.2 = \frac{1.2}{s+1}
\]

The actual loop transfer function:
\[
L(s) = P(s)C(s) = \frac{1.2}{s+1} \cdot 5 = \frac{6}{s+1}
\]

The actual sensitivity:
\[
S(s) = \frac{1}{1 + 6/(s+1)} = \frac{s+1}{s+7}
\]

The actual disturbance transfer function:
\[
G_{yv}(s) = P(s)S(s) = \frac{1.2}{s+1} \cdot \frac{s+1}{s+7} = \frac{1.2}{s+7}
\]

\textbf{Step 4: Relative error analysis}

The relative change in $G_{yv}$ can be computed. At steady state ($s \to 0$):

Nominal: $G_{yv,0}(0) = 1/6 = 0.167$

Actual: $G_{yv}(0) = 1.2/7 = 0.171$

Relative error:
\[
\frac{G_{yv}(0) - G_{yv,0}(0)}{G_{yv,0}(0)} = \frac{0.171 - 0.167}{0.167} = 0.024 = 2.4\%
\]

Using the theoretical relation $\frac{\mathrm{d}G_{yv}}{G_{yv}} = S\frac{\mathrm{d}P}{P}$:

At low frequency where $S_0(0) = 1/6$:
\[
\frac{\Delta G_{yv}}{G_{yv}} \approx S_0(0) \cdot \frac{\Delta P}{P} = \frac{1}{6} \cdot 0.2 = 0.033 = 3.3\%
\]

The small discrepancy is due to the linearization approximation.

\textbf{Key insights:}
\begin{itemize}
\item Load disturbance attenuation is relatively insensitive to process variations at frequencies where $S$ is small
\item The sensitivity function $S$ acts as a "scaling factor" for how process uncertainty affects disturbance rejection
\item Good feedback design (small $S$ at low frequencies) provides both good disturbance rejection AND insensitivity to model errors
\end{itemize}
\end{document}
